\section*{Chapitre \ref{ch:g6:pure-substances-and-mixtures} --- Corps purs et mélanges}

\begin{solution}{exo:g6:pure-substances-and-mixtures:1}
\omterm{def:g6:pure-substances-and-mixtures:pure-substance}{Corps purs}~: l'eau distillée, le bloc de fer pur, le cristal de sucre.
\omterm{def:g6:pure-substances-and-mixtures:mixture}{Mélanges}~: l'\omterm{def:g4:air-a-mixture-of-gases:air}{air}, l'eau de mer, le lait.
\end{solution}

\begin{solution}{exo:g6:pure-substances-and-mixtures:2}
Homogènes~: l'eau salée, l'eau du robinet. Hétérogènes~: le granite,
l'huile sur l'eau, le muesli.
\end{solution}

\begin{solution}{exo:g6:pure-substances-and-mixtures:3}
Une sorte particulière de substance, avec son propre nom et ses propres
propriétés~: par exemple l'eau, le sel, l'éthanol (ou encore le sucre,
l'oxygène, le fer).
\end{solution}

\begin{solution}{exo:g6:pure-substances-and-mixtures:4}
Une température d'ébullition fixe suggère un \omterm{def:g6:pure-substances-and-mixtures:pure-substance}{corps pur}~;
\qty{78}{\celsius} est la température d'ébullition de l'éthanol.
\end{solution}

\begin{solution}{exo:g6:pure-substances-and-mixtures:5}
Il contient beaucoup d'\omterm{def:g6:pure-substances-and-mixtures:species}{espèces chimiques} (eau, sucres, acides,
vitamines, arômes)~: c'est un \omterm{def:g6:pure-substances-and-mixtures:mixture}{mélange}. «~Pur~» sur la brique signifie
seulement que rien n'a été ajouté.
\end{solution}

\begin{solution}{exo:g6:pure-substances-and-mixtures:6}
Non~: un \omterm{def:g6:pure-substances-and-mixtures:pure-substance}{corps pur} bout à une température fixe. C'est un \omterm{def:g6:pure-substances-and-mixtures:mixture}{mélange},
probablement de l'eau dans laquelle quelque chose est dissous, comme de
l'eau salée.
\end{solution}

\begin{solution}{exo:g6:pure-substances-and-mixtures:7}
Un filtre pourrait séparer l'eau boueuse (les grains restent sur le
papier) et le jus avec pulpe (la pulpe reste). Il ne peut séparer ni
l'eau salée (le sel est dissous) ni l'huile de l'eau (les deux liquides
passent~; on les sépare en recueillant la couche du dessus).
\end{solution}

\begin{solution}{exo:g6:pure-substances-and-mixtures:8}
$78.7 \div 10 = \qty{7.87}{g}$ pour \qty{1}{cm^3}~: du fer.
\end{solution}

\begin{solution}{exo:g6:pure-substances-and-mixtures:9}
$\frac{357}{478} \approx 0.747$~: environ \qty{75}{\percent}.
\end{solution}

\begin{solution}{exo:g6:pure-substances-and-mixtures:10}
$2 \times 35 = \qty{70}{g}$ de sels. Pourcentage~: $\frac{35}{1000} =
\qty{3.5}{\percent}$.
\end{solution}

\begin{solution}{exo:g6:pure-substances-and-mixtures:11}
Mesurer leurs températures d'ébullition (chacune doit rester fixe
pendant l'ébullition, et les deux doivent être égales), et peser un même
volume de chacun (masses égales). Deux constantes égales, toutes deux
fixes, désignent le même \omterm{def:g6:pure-substances-and-mixtures:pure-substance}{corps pur}.
\end{solution}

\begin{solution}{exo:g6:pure-substances-and-mixtures:12}
$\frac{10}{90 + 10} = \qty{10}{\percent}$ d'éthanol. C'est un \omterm{def:g6:pure-substances-and-mixtures:mixture}{mélange}~:
il n'a pas de température d'ébullition fixe et n'apparaît pas dans la
table, qui ne contient que des \omterm{def:g6:pure-substances-and-mixtures:pure-substance}{corps purs}.
\end{solution}

\begin{solution}{pb:g6:pure-substances-and-mixtures:1}
\textbf{1.} Non. Un \omterm{def:g6:pure-substances-and-mixtures:pure-substance}{corps pur} et un \omterm{def:g6:pure-substances-and-mixtures:homogeneous}{mélange homogène} peuvent avoir
exactement le même aspect~: il faut mesurer quelque chose.

\textbf{2.} Homogène~: il est limpide et on ne voit pas ses parties.

\textbf{3.} Une constante~: une température d'ébullition qui reste fixe
pendant l'ébullition (ou la masse d'un volume connu, comparée à une
table).

\textbf{4.} A et B, dont les températures d'ébullition restent fixes,
sont des \omterm{def:g6:pure-substances-and-mixtures:pure-substance}{corps purs}. C est un \omterm{def:g6:pure-substances-and-mixtures:mixture}{mélange}.

\textbf{5.} A bout à \qty{100}{\celsius}~: c'est l'eau. B bout à
\qty{78}{\celsius}~: c'est l'éthanol.

\textbf{6.} A~: $49.8 \div 50.0 = \qty{0.996}{g}$. B~: $39.5 \div 50.0 =
\qty{0.790}{g}$. C~: $51.2 \div 50.0 = \qty{1.024}{g}$.

\textbf{7.} Oui~: environ \qty{1.0}{g} pour l'eau, \qty{0.79}{g} pour
l'éthanol.

\textbf{8.} À mesure que l'eau s'évapore, le sel reste~: le liquide
restant contient de plus en plus de sel, et il bout alors à une
température plus élevée.

\textbf{9.} Le liquide C contient un solide dissous~: c'est une
solution, un \omterm{def:g6:pure-substances-and-mixtures:mixture}{mélange}, comme son ébullition l'avait déjà montré.

\textbf{10.} $100 \times 1.024 = \qty{102.4}{g}$~; $\frac{3.5}{102.4}
\approx 0.034$~: environ \qty{3.4}{\percent} de la masse est du sel.

\textbf{11.} $\frac{35}{1000} = \qty{3.5}{\percent}$~: très proche de la
question 10.

\textbf{12.} \qty{3.5}{g} dans \qty{100}{mL}, donc $3.5 \times 10 =
\qty{35}{g}$ de sel par litre. Par sa salinité et par son ébullition, C
est l'eau de mer~; A est l'eau distillée et B l'éthanol.
\end{solution}
